\section{Experimental Methodology}\label{sec:method}
\subsection{Prediction problem and experimental stages}
Each estimator maps a feature row $\mathbf{x}_t$ to one hourly total $\widehat y_t$. One pooled model serves all target hours within a setting, rather than 24 separate lead-specific estimators. The baseline, default, tuned and ablation comparisons share eligible target timestamps within their respective settings.

The staged protocol is: L0/L1 establish seasonal baselines and default regressors on November; L2 applies default models to five temporal folds; L3 searches parameters and compares absolute/residual targets on those same folds; notebook 06 tests feature subsets and separately tunes h1; notebook 07 refits on development data and evaluates December; notebook 08 performs post-final diagnostic experiments. The stages are a final organization of the code, not a claim that feature choices were made only once. The development history records iterative feature and target choices followed by reruns of earlier stages.

\subsection{Partitions and expanding-window validation}
The fixed boundaries are 1 November and 1 December 2017. Training for L1 ends on 31 October; November is its validation month. Development combines those two periods. Test observations span 1--30 December. \Cref{tab:splits} gives actual sample counts after warm-up. The h1 development set has 23 more early January hours because of its shorter warm-up; the validation and test periods are identical.
\begin{table}[tbp]\centering\small
\caption{Chronological partitions after history-related row removal (T02). All dates are in 2017; end times are 23:00.}\label{tab:splits}
\setlength{\tabcolsep}{4pt}
\begin{tabular}{llrrr}
\toprule
Setting & Partition & Start & End & Hours \\
\midrule
h24 & Train & 01-08 23:00 & 10-31 23:00 & \num{7105} \\
h24 & Validation & 11-01 00:00 & 11-30 23:00 & \num{720} \\
h24 & Development & 01-08 23:00 & 11-30 23:00 & \num{7825} \\
h24 & Test & 12-01 00:00 & 12-30 23:00 & \num{720} \\
h1 & Train & 01-08 00:00 & 10-31 23:00 & \num{7128} \\
h1 & Validation & 11-01 00:00 & 11-30 23:00 & \num{720} \\
h1 & Development & 01-08 00:00 & 11-30 23:00 & \num{7848} \\
h1 & Test & 12-01 00:00 & 12-30 23:00 & \num{720} \\
\bottomrule\end{tabular}
\end{table}


\code{TimeSeriesSplit(n_splits=5, test_size=720)} creates five expanding training windows with 30-day validation windows\cite{tscv}. \Cref{tab:folds,fig:split} make the boundaries explicit. The last fold exactly reproduces the L1 train/November validation partition. Fold numbering starts at zero, matching the saved tables.
\begin{table}[H]\centering\small
\caption{Expanding-window h24 cross-validation (T02). Training starts on 8 January at 23:00 in every fold; validation begins at 00:00 and ends at 23:00.}\label{tab:folds}
\setlength{\tabcolsep}{4pt}
\begin{tabular}{rlrll}
\toprule
Fold & Training ends & Train hours & Validation starts & Validation ends \\
\midrule
\num{0} & 2017-07-03 & \num{4225} & 2017-07-04 & 2017-08-02 \\
\num{1} & 2017-08-02 & \num{4945} & 2017-08-03 & 2017-09-01 \\
\num{2} & 2017-09-01 & \num{5665} & 2017-09-02 & 2017-10-01 \\
\num{3} & 2017-10-01 & \num{6385} & 2017-10-02 & 2017-10-31 \\
\num{4} & 2017-10-31 & \num{7105} & 2017-11-01 & 2017-11-30 \\
\bottomrule\end{tabular}
\end{table}

\begin{figure}[tbp]\centering
\includegraphics[width=\linewidth]{F10_split_and_folds.png}
\caption{Chronological partitions and five expanding validation folds (F10). December is outside model-selection folds. The original graphic's ``untouched'' wording describes the intended partition; whole-year descriptive EDA and later test reruns qualify literal unseen-data claims.}\label{fig:split}
\end{figure}

No gap is specified, so the implementation uses zero gap. Under the midnight interpretation, completed prior-day labels can be available when the next validation day begins; load-history columns do not access later target-day loads. This assumes negligible collection and model-fitting latency. It does not establish the validity of a different issuance time, nor does it resolve unavailable future weather. The repository contains no saved comparison with a nonzero gap.

The model is fitted once per fold and evaluated with measured historical loads updating through the validation window. This differs from an open-loop forecast of all 720 hours at the fold boundary. Fold SD reflects changing regimes and training sizes. Overlapping training sets and dependent hourly errors mean it is not a standard error of five independent estimates.

\subsection{Baselines and model families}
The seasonal baselines are $\widehat y_t=y_{t-24}$ and $\widehat y_t=y_{t-168}$. MeanProfile estimates the training mean for $(a_t,d_t)$, with the training global mean as a fallback for an absent combination. The h1 persistence baseline uses $y_{t-1}$. Baselines are evaluated on the same rows as ML models, using the full constructed feature matrix where necessary.

The default model implementations are standardized ordinary least squares, an unconstrained decision tree, a 100-tree random forest, standardized RBF SVR with target scaling, and XGBRegressor. Stochastic estimators use seed 42. RF-P1 uses 30 trees and seven candidate features per split, with unrestricted depth, minimum split size two and minimum leaf size one. This imports a literature configuration into the current pipeline; it is not a reproduction of that paper's complete experiment.

In L3, the LR label denotes Ridge, and each family has an absolute and residual version. Parameters not explicitly set are inherited from the repository's pinned libraries. \Cref{sec:tuning} gives the search details. No additional forecasting family is silently introduced.

\subsection{Metrics and aggregation}
For $n$ evaluated hours with observations $y_i$, forecasts $\widehat y_i$ and evaluation-set mean $\bar y$, the implemented metrics are
\begin{align}
\mathrm{MAE}&=\frac{1}{n}\sum_{i=1}^{n}|y_i-\widehat y_i|,\\
\mathrm{RMSE}&=\sqrt{\frac{1}{n}\sum_{i=1}^{n}(y_i-\widehat y_i)^2},\\
\mathrm{MAPE}&=\frac{100}{n}\sum_{i=1}^{n}\left|\frac{y_i-\widehat y_i}{y_i}\right|,\\
R^2&=1-\frac{\sum_i(y_i-\widehat y_i)^2}{\sum_i(y_i-\bar y)^2}.
\label{eq:metrics}
\end{align}
RMSE emphasizes large errors; MAE measures average absolute deviation. MAPE weights errors relative to each actual load and is reported in percent. The custom implementation has no special zero-target treatment, but observed loads are positive. R$^2$ depends on the variance of the evaluation period, so differences across settings require care.

CV reports the arithmetic mean of five fold metrics and sample SD. Mean fold RMSE is generally different from RMSE pooled over all fold predictions, even with equal fold sizes. Likewise, mean fold R$^2$ is not pooled R$^2$. Search-result SDs use a population convention internally; this report's displayed fold SD follows the saved pandas summaries.

The baseline-relative skill is
\begin{equation}
S=1-\frac{\mathrm{RMSE}_{model}}{\mathrm{RMSE}_{Naive24}}.
\label{eq:skill}
\end{equation}
Positive skill indicates improvement; negative skill indicates worse RMSE. The L2 and ablation summaries average fold-specific skills, while the tuning comparison includes a ratio of mean RMSEs. Those summaries should not be interchanged. December skill compares metrics on the same 720 hours. Additional diagnostics include signed bias, error by hour/day, runtime and validation-to-training RMSE ratio; none replaces mean CV RMSE as the selection objective.

\subsection{Selection and evaluation limitations}
The final h24 identity is the minimum mean L3 CV RMSE among tuned variants. The final h1 identity follows its separate development search. Both select absolute-target SVR. Although notebook 07 computes a test table before printing the chosen name, the choice expression reads only development CV results. The static code therefore supports the stated selection rule, but is not a prospective registration of the development process.

Feature selection, holiday interpretation, residual-base selection and parameter optimization reuse January--November observations. No nested outer CV or correction for repeated selection is implemented. Whole-year descriptive EDA exposes the test month; progress notes also report later unchanged-pipeline reruns. The appropriate description is a separate chronological test period under a documented CV selection rule, with limits to one-shot independence. No significance test, calibration analysis or predictive interval is reported.
\evidence{\code{src/config.py}; \code{src/evaluate.py}; \code{src/models.py}; T02; notebooks 03--08; \code{docs/progress.md}.}
